% BHU PYQ -- Manually transcribed & error-corrected
\documentclass[11pt,a4paper]{article}

\usepackage[a4paper,top=2.5cm,bottom=2.5cm,left=2.8cm,right=2.8cm,headheight=14pt]{geometry}
\usepackage{amsmath,amssymb}
\usepackage[T1]{fontenc}
\usepackage[utf8]{inputenc}
\usepackage{lmodern}
\usepackage{microtype}
\usepackage{fancyhdr}
\usepackage{enumitem}
\usepackage{hyperref}
\usepackage{lastpage}
\usepackage{parskip}

\hypersetup{colorlinks=true,urlcolor=black,linkcolor=black,
  pdftitle={CHB-201: Inorganic & Physical Chemistry-I -- BHU PYQ},pdfauthor={Banaras Hindu University}}

\pagestyle{fancy}\fancyhf{}
\renewcommand{\headrulewidth}{0.4pt}\renewcommand{\footrulewidth}{0.4pt}
\fancyhead[L]{\small   \textit{Banaras Hindu University}}
\fancyhead[R]{\small   \textit{CHB-201: Inorganic \& Physical Chemistry-I}}
\fancyfoot[C]{\small Page  \thepage\ of \pageref{LastPage}}


\newlist{parts}{enumerate}{2}
\setlist[parts,1]{label=(\alph*),leftmargin=*,itemsep=5pt,topsep=3pt}
\setlist[parts,2]{label=(\roman*),leftmargin=*,itemsep=3pt,topsep=2pt}

\newcommand{\pts}[1]{\hfill   \textit{[#1]}}


\begin{document}


\begin{center}
  {\large\bfseries BANARAS HINDU UNIVERSITY}\\[2pt]
  {\normalsize B.Sc. (Hons.) Semester II Examination, 2013-14}\\[6pt]
  \rule{0.8\linewidth}{0.5pt}\\[6pt]
  {\large\bfseries Paper No. CHB-201: Inorganic \& Physical Chemistry-I}\\[4pt]
  {\normalsize Time: 3 Hours\hfill Maximum Marks: 70}
\end{center}
\rule{\linewidth}{0.4pt}
\small



\noindent   \textit{Instructions: Answer five questions in total, including Question~1 from each section which is compulsory. Use separate answer books for Section~A and Section~B.}

\normalsize
\bigskip


\begin{center}
  {\large\bfseries SECTION A}\\[2pt]
  {\normalsize (Inorganic Chemistry-I -- Marks: 35)}
\end{center}
\rule{0.4\linewidth}{0.2pt}
\smallskip




\noindent   \textbf{Question 1.} Answer all parts of the following: \pts{5 $ \times$ 2 = 10}

\begin{parts}
  \item Give a valid reason for the high solubility of $ \text{BeF}_2$ in water, while all other alkaline earth fluorides are insoluble.
  \item Write the structures of $ \text{BeCl}_2$ in both the vapour phase and solid state.
  \item Write the molecular formula of 'Borax' and draw its structure.
  \item Write the formula of Beryl. To which class of silicates does it belong?
  \item Name the anhydride of $ \text{HNO}_3$ and draw its gas-phase structure.
\end{parts}

\medskip\rule{\linewidth}{0.2pt}




\noindent   \textbf{Question 2.}

\begin{parts}
  \item State the Kapustinskii equation for lattice energy and calculate the lattice energy ($\Delta H_L$) of $ \text{NaCl}$ using a Born-Haber cycle with the following data (all values in $ \text{kJ/mol}$):
    \[ \Delta H_f^\circ = -411.2, \quad \Delta H_{ \text{sub}} = 108.4, \quad  \text{IE} = 495.4, \quad \Delta H_{ \text{diss}}( \text{Cl}_2) = 241.8, \quad  \text{EA}( \text{Cl}) = -348.6 \] \pts{6}
  \item Define 'electrochemical series'. Answer the following, giving reasons, based on the standard reduction potentials ($E^\circ$ in volts):
    \[  \text{Zn}^{2+}/ \text{Zn} = -0.76, \quad  \text{Fe}^{2+}/ \text{Fe} = -0.44, \quad  \text{Cu}^{2+}/ \text{Cu} = +0.34, \quad  \text{Ag}^+/ \text{Ag} = +0.80 \]
  
\begin{parts}
    \item Can we store copper sulphate solution in a zinc container?
    \item Which is the strongest reducing agent among these metals?
  \end{parts} \pts{6}
\end{parts}

\medskip\rule{\linewidth}{0.2pt}




\noindent   \textbf{Question 3.}

\begin{parts}
  \item Answer the following with respect to diborane ($ \text{B}_2 \text{H}_6$):
  
\begin{parts}
    \item Name the experimental or physical property which confirmed that rotation between the two boron atoms is hindered.
    \item Give two experimental evidences which support that four terminal hydrogen atoms are in a different environment from the two bridging hydrogen atoms.
  \end{parts} \pts{6}
  \item Draw the structures of hydrazine ($ \text{N}_2 \text{H}_4$) and hydroxylamine ($ \text{NH}_2 \text{OH}$). Explain why hydrazine is a weaker base than ammonia. \pts{6}
\end{parts}

\medskip\rule{\linewidth}{0.2pt}




\noindent   \textbf{Question 4.}

\begin{parts}
  \item Show hybrid orbitals, draw structures, and mention the shapes of the following: (i) $ \text{XeO}_2 \text{F}_2$ and (ii) $ \text{XeO}_3 \text{F}_2$. \pts{6}
  \item How do you name the $[ \text{XeO}_6]^{4-}$ ion? Show hybrid orbitals, mention the number of $\sigma$- and $\pi$-bonds, and draw the structure of the ion. \pts{6}
\end{parts}


\newpage

\begin{center}
  {\large\bfseries SECTION B}\\[2pt]
  {\normalsize (Physical Chemistry-I -- Marks: 35)}
\end{center}
\rule{0.4\linewidth}{0.2pt}
\smallskip




\noindent   \textbf{Question 5.} Answer all parts of the following: \pts{5 $ \times$ 2 = 10}

\begin{parts}
  \item What is the influence of temperature on the surface tension and viscosity of a liquid?
  \item Define critical temperature ($T_c$), Boyle temperature ($T_B$), and inversion temperature ($T_i$).
  \item Heat change ($q$) is a path-dependent thermodynamic quantity. State the conditions under which it becomes a state function.
  \item Show that for a reversible adiabatic expansion of an ideal gas, the work done is given by $W = C_v(T_1 - T_2)$.
  \item Define the term 'mean free path' of a gas molecule.
\end{parts}

\medskip\rule{\linewidth}{0.2pt}




\noindent   \textbf{Question 6.}

\begin{parts}
  \item Derive Kirchhoff's equations for the variation of enthalpy of reaction with temperature. \pts{6}
  \item Show that the entropy change in an isothermal reversible process for the universe is zero, while it is positive in a spontaneous (irreversible) process. \pts{6}
\end{parts}

\medskip\rule{\linewidth}{0.2pt}




\noindent   \textbf{Question 7.}

\begin{parts}
  \item For reversible isothermal and adiabatic expansion of an ideal gas, demonstrate how thermodynamic quantities ($w$, $q$, $\Delta U$, and $\Delta H$) can be estimated. \pts{6}
  \item Using the kinetic gas equation, deduce Avogadro's law and Graham's law of diffusion. \pts{6}
\end{parts}

\medskip\rule{\linewidth}{0.2pt}




\noindent   \textbf{Question 8.}

\begin{parts}
  \item Define the coefficient of viscosity of a liquid. Discuss the Ostwald viscometer method for the determination of the viscosity of a liquid. \pts{6}
  \item Write the van der Waals equation of state for a real gas. Define critical constants ($P_c$, $V_c$, $T_c$) in terms of van der Waals constants ($a$ and $b$). \pts{6}
\end{parts}

\vfill
\rule{\linewidth}{0.4pt}

\begin{center}\small
  \href{https://bhu-exam.vercel.app/}{Click here to visit our website}\\[4pt]
  \href{https://me-aryan.vercel.app/}{Click here to visit our contributor}\\[4pt]
  \href{https://quashan-me.vercel.app/}{Click here to visit our contributor}
\end{center}

\end{document}
